% Рисунок с доски: характеристический треугольник с вершиной (x_0, t_0).
% Боковые стороны -- характеристики |x - x_0| = a(t_0 - t).
\begin{figure}[htbp]
\centering
\begin{tikzpicture}[>={Latex[length=4pt]}]
  \draw[->] (-0.4,0) -- (6.0,0) node[anchor=north east] {$x$};
  \draw[->] (0,-0.4) -- (0,3.4) node[anchor=south east] {$t$};

  \fill[blue!10] (1.0,0) -- (5.0,0) -- (3.0,2.6) -- cycle;
  \draw[thick, blue!60!black] (1.0,0) -- (5.0,0) -- (3.0,2.6) -- cycle;

  \fill (3.0,2.6) circle (1.6pt);
  \node[anchor=south, font=\footnotesize] at (3.0,2.68) {$(x_0,t_0)$};
  \node[anchor=west, font=\footnotesize] at (4.15,1.65) {$|x - x_0| = a(t_0 - t)$};

  % промежуточный срез t = tau
  \draw[dashed, gray] (0,1.3) -- (2.0,1.3);
  \draw[dashed, gray] (4.0,1.3) -- (4.6,1.3);
  \draw[very thick, red!60!black] (2.0,1.3) -- (4.0,1.3);
  \node[anchor=east, font=\footnotesize] at (-0.08,1.3) {$\tau$};
  \node[anchor=east, font=\footnotesize] at (-0.08,2.6) {$t_0$};
  \draw[dashed, gray] (0,2.6) -- (3.0,2.6);

  \node[anchor=north, font=\footnotesize] at (3.0,-0.05) {$x_0$};
  \node[anchor=north, font=\footnotesize] at (1.0,-0.05) {$x_0 - a t_0$};
  \node[anchor=north, font=\footnotesize] at (5.0,-0.05) {$x_0 + a t_0$};
  \node[blue!60!black, font=\footnotesize] at (3.0,0.6) {$D$};
\end{tikzpicture}
\caption{Характеристический треугольник $D$ с вершиной $(x_0,t_0)$: прообраз
области $D$ из рис.~\ref{fig:char-xieta} при возврате к переменным $(x,t)$.
Горизонтальный срез на уровне $t = \tau$ — отрезок
$\big[x_0 - a(t_0-\tau),\ x_0 + a(t_0-\tau)\big]$.}
\label{fig:char-triangle}
\end{figure}
